\title{Search Smooths Discontinuities in\\ Discrete/Continuous Problems}

\author{
Christopher D.\ Carroll\thanks{Carroll: Department of Economics, Johns Hopkins University,
\url{http://econ.jhu.edu/people/ccarroll/}, \texttt{ccarroll@jhu.edu}.}\\
JHU
\and
Fedor Iskhakov\thanks{Iskhakov: Stony Brook University.}\\
Stony Brook University
\and
Shujaat Khan\thanks{Khan: IMF.}\\
IMF
\and
Adrian Monninger\thanks{Monninger: Department of Economics, Texas A\&M University, \url{https://sites.google.com/view/adrianmonninger}, \texttt{amonninger@tamu.edu}.}\\
Texas A\&M
}

\date{\today}

\maketitle

\begin{abstract}
% Previous abstract:
% When used to solve models with both discrete and continuous choices, the Endogenous Gridpoint Method (EGM), in addition to finding the optimal solution, can also generate points in the solution that are suboptimal. This paper shows that incorporating search into such problems can smooth out the discontinuities in the decision rules, which are responsible for producing the suboptimal points. Combining search with price dispersion leads to ample smoothing which can eliminate the suboptimal points entirely. In addition to providing greater control over the smoothness of the problem, search offers the added benefit of rationalizing the smoothness into the agent's choice problem.
Discrete-continuous problems solved by the endogenous gridpoint method produce kinks in the value function and jumps in the decision rule. Those jumps make the Euler equation necessary but not sufficient, so the method returns suboptimal points. Extreme Value Type~I taste shocks remove the jumps \citep{Iskhakov2017-jb}. This paper shows that directed search together with a rejectable price draw can match that smoothness, using objects that are part of the agent's problem. The agent chooses effort in the unit interval at a convex utility cost, and effort raises the probability of receiving an option to switch. Conditional on receiving the option, the agent draws a price and may reject it. A retirement model illustrates the mechanism. In a keep-versus-adjust model of durables, taste shocks, price dispersion, and search all smooth the $(S,s)$ band, and their steady-state effects on the adjustment probability have opposite signs. Those signs identify the search cost and the price dispersion separately. Matching impulse responses of consumption and of the adjustment probability remains open.

\medskip
% Previous keywords:
% \noindent\textbf{Keywords:} Consumption, Computational Economics, Discrete Continuous Choice, Endogenous Gridpoint Method, Search
\noindent\textbf{Keywords:} Consumption, Computational Economics, Discrete-Continuous Choice, Endogenous Gridpoint Method, Search, Price Dispersion

\noindent\textbf{JEL codes:} C13, C63

\medskip
\noindent\textbf{Repo:} \url{https://github.com/AMonninger/DCEGM_Search}\\
\noindent\textbf{PDF:} \url{https://github.com/AMonninger/DCEGM_Search/Text/dcegmSearch.pdf}
\end{abstract}

\bigskip
\noindent\footnotesize The computational results in this paper were constructed using tools in the Econ-ARK/HARK toolkit. The toolkit can be cited by its digital object identifier 10.5281/zenodo.1001068, as is done in the paper's own references as \citet{Carroll2017-hk}.
Thanks to Matthew White for initial discussions of this subject.
\normalsize
